\chapter*{Introduction}
\markleft{Introduction}
\addcontentsline{toc}{chapter}{Introduction}

Dans un contexte où les services numériques sont sollicités en permanence, la capacité d’une entreprise
à détecter rapidement les incidents et à mobiliser les bonnes personnes constitue un enjeu opérationnel
majeur. Cela est particulièrement vrai pour une plateforme de jeu en ligne comme Winamax, où la moindre
indisponibilité se traduit immédiatement par une dégradation de l’expérience de millions de joueurs, une
perte de revenus et une atteinte potentielle à la réputation de l’entreprise. Le système d’alerting, c’est-à-
dire l’ensemble des mécanismes qui détectent une anomalie et notifient les équipes d’astreinte, est au cœur
de cette réactivité.

Ce rapport présente les travaux menés durant mon stage de fin d’études chez Winamax, au sein de
l’équipe Infrastructure, autour de la refonte de ce système d’alerting. Le point de départ du sujet est un
constat concret et partagé par les équipes : le système d’alertes existant produit un volume important de
notifications, dont une partie significative reste peu pertinente pour l’astreinte. De nombreuses alarmes
sont en effet construites autour de métriques de causes — par exemple une consommation de CPU
ou de mémoire élevée — plutôt qu’autour de conséquences directement corrélées à un impact utilisateur.
Cette situation augmente la charge cognitive des équipes, provoque des réveils inutiles et peut retarder
l’identification du service réellement en difficulté.

L’objectif du stage est donc de faire évoluer cette approche vers un modèle de criticité plus lisible et
plus actionnable, permettant de prioriser les interventions selon l’impact réel des incidents. Au-delà du
simple paramétrage technique, le sujet implique une véritable démarche d’ingénierie : analyse de l’existant,
conception d’un modèle d’observabilité, choix d’une architecture cible, développement en Infrastructure
as Code, et réflexion sur la gouvernance et la validation continue de la qualité des alertes.

Le présent rapport est structuré en sept chapitres. Le premier présente l’entreprise Winamax, son
organisation technique et son contexte technologique. Le deuxième détaille le sujet, les motivations et les
objectifs du stage. Le troisième analyse le système d’alerting actuel et ses limites. Le quatrième décrit la
conception de la nouvelle solution, fondée sur les méthodes RED et USE et sur des niveaux de criticité
explicites. Le cinquième expose la mise en œuvre technique de cette solution. Le sixième revient sur
une expérience transverse marquante : mes deux semaines de « fast », dédiées au support quotidien des
équipes de développement. Enfin, le septième chapitre dresse le bilan du travail, les gains attendus et les
perspectives d’évolution.
